\scanpage{45}
% archive sheet number: 44
D'après la théorie d'Artin-Schreier-Witt, on a
\[
H^1(A,\mathbb Z/p)\simeq H^1(A,\underline O_A)_{\deleted{\unsure{p|s}}}^{(p)}
\qquad \text{éléments de } H^1(A,\underline O_A)\text{ tels que } Fx=x
\]
\[
H^1(A,\mathbb Z/p^{\unsure{n|m}})\simeq
H^1(\deleted{\unsure{O|A}}A,\underline O_A(W_n))^{(p)}
\qquad \text{--- }H^1(A,\underline O_A(W_n))\text{ tels que }Fx=x
\]
\[
\boxed{\operatorname{Hom}(\pi_1^p(A),\mathbb Z_p)\simeq H^1(A,\deleted{\unsure{\mathbb Q_p|\underline O_p}}\mathbb Z_p)
\simeq H^1(A,W)_{\deleted{\unsure{s|p}}}^{(p)}}
\]
\deleted{\unsure{idem|idée}} \inserted{\unsure{dont|d'où}} \deleted{\unsure{on}} on déduit un \deleted{accouplement} isomorphisme
\[
\boxed{\deleted{\pi_1^p(A)\times H^1(A,W)_s\longrightarrow\mathbb Z_p}}
\]
qui est une dualité \deleted{\unsure{de|(?)}} de $\mathbb Z_p$-modules libres de type fini
\[
\boxed{\deleted{\operatorname{Hom}_{\unsure{\mathrm{cont}|\mathrm{ant}}}(\pi_1^p(A),\Lambda)\simeq H^1(A,W)_s}}
\]

D'autre part, utilisant encore le fait que tout revêtement de $A$ est une variété abélienne, on trouve aussi
\[
\boxed{\pi_1^p(A)\simeq T_p(A)}
\]
d'où en les \unsure{comparant|composant}
\[
\boxed{T_p(A)\times H^1(A,W)^{(p)}\longrightarrow\mathbb Z_p}
\]
qui est une dualité de $\mathbb Z_p$-modules libres de type fini.

\deleted{qu'on peut aussi écrire} \unsure{dont|donc|d'où|dans} on déduit
\[
\boxed{\operatorname{Hom}_{\unsure{\mathrm{cont}|\mathrm{ant}}}(T_p(A),\deleted{\unsure{\mathbb Z|\Lambda}}W_{\infty}(k))\simeq H^1(A,W)_s}
\]
